\section{Experimental design, inference and computational budget}\label{sec:section-05}

The controlled study completed 848 registered trials and 1,696 estimator fits, in addition to 4 development runtime-pilot fits. Its base evaluation spans 252 outer dates and 5,544 unique base instrument-date decisions. The registered families contain 17 predictive contrasts and 8 economic contrasts. Repeated model views and seed runs are retained, but never counted as additional independent market observations.

\noindent Table~\ref{tab:4} reports frozen controlled training design.

\begin{table}[!htbp]\centering\footnotesize\setlength{\tabcolsep}{3pt}\renewcommand{\arraystretch}{1.14}
\caption{Frozen controlled training design}\label{tab:4}
\begin{tabular}{@{}>{\raggedright\arraybackslash}p{\dimexpr 0.34646\textwidth-4.800pt\relax}>{\raggedleft\arraybackslash}p{\dimexpr 0.15354\textwidth-4.800pt\relax}>{\raggedleft\arraybackslash}p{\dimexpr 0.15354\textwidth-4.800pt\relax}>{\raggedleft\arraybackslash}p{\dimexpr 0.17323\textwidth-4.800pt\relax}>{\raggedleft\arraybackslash}p{\dimexpr 0.17323\textwidth-4.800pt\relax}@{}}
\toprule
\textbf{Variant} & \textbf{Trials} & \textbf{Targets} & \textbf{Outer folds} & \textbf{Seeds} \\
\midrule
base & 136 & 22 & 4 & 42,137,2026 \\
without\_\allowbreak{}momentum & 88 & 22 & 4 & 42 \\
without\_\allowbreak{}trend & 88 & 22 & 4 & 42 \\
without\_\allowbreak{}risk & 88 & 22 & 4 & 42 \\
without\_\allowbreak{}context\_\allowbreak{}volume & 88 & 22 & 4 & 42 \\
fraction\_\allowbreak{}0.5 & 88 & 22 & 4 & 42 \\
fraction\_\allowbreak{}0.75 & 88 & 22 & 4 & 42 \\
history\_\allowbreak{}504 & 88 & 22 & 4 & 42 \\
cadence\_\allowbreak{}21 & 72 & 6 & 4 & 42 \\
transfer & 24 & 6 & 4 & 42 \\
\bottomrule\end{tabular}\par\smallskip
\raggedright\textit{Source:} \path{data/research-v2/experiments/design.json}.
\end{table}


The primary predictive comparison averages each day's candidate-minus-prior log-loss difference across symbols, then resamples contiguous date blocks. Keeping the cross section intact avoids treating a market-wide error as many independent surprises. Circular wrapping allows every date to start a block, but short samples provide limited distinct market episodes. Block-length sensitivity examines a label-horizon-scale block, the registered primary block and a longer block. Confidence intervals describe the behavior of this historical sample under the resampling assumptions; they do not make the consumed audit confirmatory.

Economic comparisons use identically indexed daily net returns. Rebased charts make date-range returns visually comparable but do not reset account holdings. Where a fresh-capital replay is reported, it explicitly begins a new account and is separately labeled. CAGR compounds the observed return over its annualized duration. The new portfolio diagnostics use a zero risk-free rate for Sharpe and a zero minimum acceptable return for Sortino, consistent with zero-yield cash. The sealed release separately reports Sharpe relative to a fixed annual hurdle; those differently defined values are never interchanged. Sortino uses downside deviation relative to its disclosed target. Undefined zero-variance statistics remain undefined rather than being converted into evidence for a strategy.

Hypothesis families come from registered contrasts. Holm-adjusted values apply separately to the specified predictive and economic families. Additional class, symbol and regime slices are descriptive unless entered into a distinct family before computation. This avoids presenting a favorable post hoc slice as a new independent experiment. Effect sizes and intervals take precedence over binary significance language. Approximate minimum-detectable-effect calculations use the estimated sampling variability; they are design aids, not guarantees of future power.

Execution robustness perturbs clearly named assumptions. Cost changes can retain decisions while changing financing and fills. Delaying execution or restricting fill size may alter exposure and subsequent account behavior. Missing-bar tests are admission/recovery tests and must fail before decisions rather than interpolate a fictional observation. Risk-parameter changes create a different target policy and are classified accordingly. The table records which components remain frozen in each experiment so that a changed return is not assigned to the wrong cause.

Runtime evidence is separate from correctness evidence. Measurements disclose hardware, interpreter, thread limits, sample counts and cold/warm conditions. Quantiles from small samples are descriptive order statistics, not reliable tail service-level estimates. Compute budgets are frozen after a development-only pilot, attempted trials retain status and errors, and cache keys include data, source, configuration and fold boundaries. Exact repetition from saved source is a numerical reproduction claim; operation across a genuinely new market period is a different claim.
